\chapter{Confronting the Dual-Scale Hypothesis with Data: Pre-Registration and Four Verdicts}\label{ch:verdicts}

A physical hypothesis earns its name only if some measurement could prove it wrong. The
dual-scale programme of this book began with an attractive idea: treat a microscopic length and
a macroscopic length as a T-dual pair, as string theory treats $R$ and $\ap/R$ in
\cref{ch:tduality}. \Cref{ch:cosmology} and \cref{ch:open} showed what the formal part of that
idea amounts to. This chapter asks the question a student should ask of any proposal: \emph{what
does it predict that we can measure, and does the measurement agree?} Answering it honestly
takes more than arithmetic. A test has to be written down, with its decision rule, \emph{before}
the verdict is computed, and anything the author already knew about the data has to be
declared. We first explain what such a pre-registered test is and why each of its ingredients
matters. We then list everything the programme can derive without a compactification that has
never been constructed; there turn out to be exactly four observables. We confront each with
data, derive every step on the page, and end with the single cause behind all four verdicts and
what it teaches. The Lean modules attached to the chapter check the algebra and the arithmetic
of each verdict; none of them checks that the identifications with physics are right.

Conventions: $\hbar=c=1$ unless units are written out; $\ell_P$ is the Planck length, so
$G=\ell_P^2$; $c/H_0$ is the Hubble radius today. Numbers for $\ell_P$ and $c/H_0$ are those
fixed in \lean{DualScaleCosmology/SelfDualCutoff.lean}:
$\ell_P=1.616255\times10^{-35}$\,m (CODATA 2018) and $c/H_0=1.3672\times10^{26}$\,m (Planck 2018,
$H_0=67.66$\,km/s/Mpc), both taken from \texttt{astropy} and not checked against a source in
this book's library.

\section{What makes a test}\label{sec:verdicts:test}

\subsection{Freeze first, compute second}\index{pre-registration}

A \emph{pre-registered prediction} is a statement written and time-stamped before the comparison
with data. In this programme the time stamp is a git tag: the prediction document is committed,
tagged (for example \texttt{stream6-p1-frozen}), and pushed, and only then is the verdict
computed. The document contains three things:
\begin{enumerate}[leftmargin=*,itemsep=1pt]
\item the \emph{statement}, with every assumption labelled and every input number fixed;
\item the \emph{tests}, each naming a pinned data source and the inequality the prediction must
 satisfy to pass;
\item the \emph{decision rule}: which pattern of passed and failed tests counts as
 ``consistent'', ``excluded'' or ``in tension''.
\end{enumerate}
Once the tag exists, the document is never edited. A mistake found later is recorded in a
separate errata file (for this chapter's documents, \texttt{docs/STREAM67\_ERRATA.md}, which
corrects two line numbers and nothing else); a changed prediction is a \emph{new} prediction
with a new name.

\subsection{Disclosure: retrodiction and forward prediction}\index{retrodiction}

The authors of a prediction usually know some of the data already. If they do, the freeze
cannot make the test blind; it can only stop them from adjusting the prediction afterwards. So
each frozen document says what the authors knew. A comparison with data that were known before
the freeze is a \emph{retrodiction}; a comparison with data that did not yet exist, or had not
been looked up, is a \emph{forward prediction}. Both are useful. A retrodiction with a recorded
outcome closes a question and prevents it from being quietly reopened; a forward prediction is
the only kind that can surprise its authors.

\subsection{Is the observable empty?}\index{vacuity}

Before freezing, list the inputs of the predicted number under three headings: \emph{measured}
(taken from an experiment), \emph{derived} (proved from the assumptions, here in Lean), and
\emph{assumed} (an identification or a choice). If the assumed list is not empty, the prediction
is Tier C and must say so. A second failure mode is subtler: a Lean theorem can certify
arithmetic on numbers that were typed in by hand, and prose around it can then call the result
``consistent with observation''. The kernel checks the arithmetic, not the choice of the numbers.
The rule of this chapter is therefore that every number entering a test is either a pinned
measurement or the output of a stated derivation.

\subsection{Order-one factors must be derived}

Many physical identifications carry a factor of order one: a radius or a circumference, a
$2\pi$, a coupling of the lightest mode. When a prediction misses a bound by a factor of order
one, it is tempting to pick the convention that makes it pass. That is a fit, not a test.

\begin{pitfallbox}[Choosing the convention after reading the bound]
In \cref{sec:verdicts:p1} the predicted radius is $s\approx47\,\mu$m and the laboratory bound is
$30\,\mu$m. Writing $R=s/2\pi\approx7.5\,\mu$m instead would pass the bound, and would even land
inside the window $0.1$--$10\,\mu$m of a published ``dark dimension'' proposal
\source{MVV \S3, ll.~424--427}. Choosing $2\pi$ \emph{after} seeing the bound turns a failed test
into a successful fit. The factor must come from the theory, before the comparison; in
\cref{sec:verdicts:p1} it does, and it is $1$.
\end{pitfallbox}

\section{What the programme can predict}\label{sec:verdicts:inventory}

The question to ask first is: which results of the programme are numbers \emph{with units},
obtained by a derivation that does not pass through an explicit compactification (flux quanta,
brane wrapping numbers, moduli stabilization), since no such compactification has been
constructed for this geometry? Going through the programme:
\begin{center}\small
\begin{tabular}{@{}p{0.34\textwidth}p{0.26\textwidth}>{\raggedright\arraybackslash}p{0.32\textwidth}@{}}
\toprule
Result & Number with units? & Observable \\
\midrule
Lattices, $\Odd{d}$, double field theory, the bound $R+\ap/R\ge2\sqrt{\ap}$ (\cref{ch:dualscale}) & no & none \\
Moonshine and dyon computations (\cref{ch:moonshinecomputed,ch:dyons}) & no (integers) & none: they concern $\mathcal N=4$ string theory \\
Self-dual length $s=\sqrt{\ell_P\,c/H_0}$ & yes & P1: radius of an extra dimension \\
CKN saturation with infrared length $c/H(t)$ & yes & C-A: deceleration parameter \\
Scale-factor duality realized dynamically & with dynamics & C-B: relic gravitational waves \\
F-strings at $\ap=s^2$ & yes (dimensionless $G\mu$) & C-C: string tension \\
\bottomrule
\end{tabular}
\end{center}
All four observables rest on one identification (\TC): the string scale is the geometric mean
of the Planck length and the Hubble radius,
\begin{equation}\label{eq:verdicts:ident}
 \ap=\ell_P\cdot\frac{c}{H_0},\qquad s=\sqrt{\ap}.
\end{equation}

\section{P1: the self-dual length as an extra dimension}\label{sec:verdicts:p1}

\subsection{The prediction}

The self-dual radius of a circle is the fixed point of $R\mapsto\ap/R$ (\cref{ch:tduality}). With
\eqref{eq:verdicts:ident}, P1 identifies it with the radius $R$ of \emph{one} large extra
dimension, with the Kaluza--Klein tower starting at mass $m\sim1/R$ (the convention
$m^{-1}\sim l$ of \source{MVV \S3, l.~327}).

\begin{example}[Computing $s$]\label{ex:verdicts:s}
With the inputs above,
\begin{align*}
 s^2&=\ell_P\cdot\frac{c}{H_0}
 =1.616255\times10^{-35}\,\text{m}\times1.3672\times10^{26}\,\text{m}
 =2.2097\times10^{-9}\,\text{m}^2,\\
 s&=4.70\times10^{-5}\,\text{m}\approx47\,\mu\text{m}.
\end{align*}
The tests compare squares, $s^2$ against $(30\,\mu\text{m})^2=9\times10^{-10}\,\text{m}^2$, so no
square root is ever taken in the Lean statements.
\end{example}

\subsection{The data and the decision rule}

Three bounds were frozen with P1 (tag \texttt{stream6-p1-frozen}):
\begin{itemize}[leftmargin=*,itemsep=1pt]
\item T1, the E\"ot-Wash torsion balance: the largest extra dimension ``must have a toroidal
 radius less than $30\,\mu$m'' \source{EW, l.~289};
\item T2, the same experiment: any gravitational-strength Yukawa interaction has range
 $\lambda<38.6\,\mu$m \source{EW, ll.~283--285}, applied with $\lambda=R$ for the lightest mode
 (an identification, \TC);
\item T3, heating of old neutron stars by trapped Kaluza--Klein gravitons: for one extra
 dimension $m^{-1}\sim l<44\,\mu$m \source{MVV \S3, ll.~324--327}.
\end{itemize}
Decision rule: \emph{consistent} if all pass, \emph{excluded} if T1 fails, \emph{in tension} if
only T2 or T3 fails. The authors knew these bounds before the freeze; P1 is a retrodiction.

\subsection{The verdict}

Since $s^2=2.21\times10^{-9}\,\text{m}^2$ exceeds $(30\,\mu\text{m})^2$, $(38.6\,\mu\text{m})^2
=1.49\times10^{-9}\,\text{m}^2$ and $(44\,\mu\text{m})^2=1.94\times10^{-9}\,\text{m}^2$, all three
tests fail, and by the rule \textbf{P1 is excluded}. The margins are not rounding effects:
$s>1.5\times30\,\mu$m and $s>1.06\times44\,\mu$m. To pass T1 one would need $R=\kappa s$ with
$\kappa<30/47.0=0.64$.

\subsection{The order-one factor is fixed by T-duality}

Is $\kappa$ free? The programme's own duality answers no. The map $R\mapsto\ap/R=s^2/R$ has a
fixed point where $R=s^2/R$, i.e.\ $R^2=s^2$; for $R>0$ and $s>0$ the only solution is $R=s$.
At that radius the momentum scale $1/R$ equals the winding scale $R/\ap$:
\[
 \frac1R\Big|_{R=s}=\frac1s=\frac{s}{s^2}=\frac{R}{\ap}\Big|_{R=s},
\]
which is what makes the self-dual radius special (\cref{ch:tduality}). With the standard
Kaluza--Klein normalization $m_n=n/R$ the radius is therefore $s$ itself: $\kappa=1$. If both
circles of $T^2$ were large (two extra dimensions), the neutron-star bound becomes
$l<1.6\times10^{-4}\,\mu$m \source{MVV \S3, l.~329}, missed by a factor of about $3\times10^5$.

\begin{remark}
\Cref{ch:open} called this reading ``disfavored, not excluded'', because the order-one
factor was still open. The derivation above closes it; the pre-registered rule then classifies
P1 as excluded. The two statements are consistent: they differ in what was known about
$\kappa$.
\end{remark}

\section{C-A: holographic dark energy with infrared length \texorpdfstring{$c/H$}{c/H}}\label{sec:verdicts:ca}

\subsection{The Cohen--Kaplan--Nelson bound and its saturation}\index{CKN bound}

Cohen, Kaplan and Nelson (CKN) argued that an effective field theory with ultraviolet cutoff
$\Lambda$ in a box of size $L$ should not contain states inside their own Schwarzschild radius,
which requires \source{CKN, eq.~(2), ll.~75--78}
\begin{equation}\label{eq:verdicts:ckn}
 L^3\Lambda^4\lesssim L\,M_P^2 .
\end{equation}
The left side is the maximal energy $\Lambda^4L^3$ in the box, the right side the mass of a black
hole of size $L$. Candidate C-A assumes (\TC) that the vacuum energy saturates this bound at all
times with the programme's macroscopic length as infrared length, $L=c/H(t)$. Then
$\rho_\Lambda\sim\Lambda^4\sim M_P^2/L^2\propto H^2$, which we write as
\begin{equation}\label{eq:verdicts:rhoL}
 \rho_\Lambda(t)=\varepsilon\,\frac{3H(t)^2}{8\pi G},\qquad 0<\varepsilon<1,
\end{equation}
with a constant $\varepsilon$. We also assume (\TC) a flat Friedmann universe with pressureless
matter that is separately conserved, $\rho_m\propto a^{-3}$, and neglect radiation.

\subsection{The derivation, step by step}

Write $k=8\pi G/3$. The Friedmann equation is $H^2=k(\rho_m+\rho_\Lambda)$. Substituting
\eqref{eq:verdicts:rhoL}:
\[
 H^2=k\rho_m+\varepsilon H^2\quad\Longrightarrow\quad \rho_m=(1-\varepsilon)\frac{H^2}{k}.
\]
The matter fraction $\Omega_m=k\rho_m/H^2=1-\varepsilon$ is constant. With
$\rho_m=\rho_{m0}a^{-3}$,
\[
 H^2=\frac{k\rho_{m0}}{1-\varepsilon}\,a^{-3}\equiv\frac{C}{a^3},\qquad C>0 .
\]
The deceleration parameter is $q=-\ddot aa/\dot a^2$. Using $\dot a=aH$,
$\ddot a=\dot aH+a\dot H$, so $q=-1-\dot H/H^2$, and with $\dot H=\dot a\,\dd H/\dd a=aH\,\dd H/\dd a$,
\begin{equation}\label{eq:verdicts:q}
 q=-1-\frac{\dd\ln H}{\dd\ln a}=-1-\frac{a\,E'(a)}{2E(a)},\qquad E(a)=H(a)^2 .
\end{equation}
For $E=C/a^3$ we get $aE'/E=-3$, hence
\[
 q=-1+\tfrac32=\tfrac12\qquad\text{at every epoch and for every }\varepsilon .
\]
The dark component \eqref{eq:verdicts:rhoL} dilutes exactly like matter; the universe behaves as
if it contained dust only, $w_{\rm eff}=0$, and it decelerates.

\begin{physicsbox}[Why a Hubble-scale cutoff cannot accelerate]
If the dark energy is always a fixed fraction of $3H^2/8\pi G$, it can never come to dominate:
its share of the budget is frozen at $\varepsilon$. Acceleration needs a component whose density
falls more slowly than $H^2$ does in a matter era; tying the density to $H^2$ forbids exactly
that. This is the classic objection to Hubble-scale holographic dark energy, and it is why the
authors knew the outcome before the freeze.
\end{physicsbox}

\subsection{The verdict}

The frozen test (tag \texttt{stream7-ca-frozen}) compares with the present deceleration parameter
inferred from Planck 2018 in flat $\Lambda$CDM, $\Omega_\Lambda=0.68885$ (the value in
\lean{DarkEnergyScale.}\allowbreak\lean{omegaLambda}; not checked against a source in this book's library),
$\Omega_m\approx1-\Omega_\Lambda$:
\[
 q_0=\frac{\Omega_m}{2}-\Omega_\Lambda=\frac{0.31115}{2}-0.68885=-0.533<0 .
\]
C-A requires $q_0=+1/2$: \textbf{C-A is excluded}. Two caveats were recorded in advance: the
number $q_0$ comes from a $\Lambda$CDM fit, so it is model-dependent; and the test is a
retrodiction with a foreseeable outcome.

\section{C-B: relic gravitons from a pre-big-bang phase}\label{sec:verdicts:cb}

\subsection{The spectrum}\index{pre-big bang}

Scale-factor duality (\cref{ch:cosmology}) suggests a phase before the big bang in which the
universe contracts in the Einstein frame while the dilaton grows. Tensor perturbations amplified
during such a phase leave a relic background of gravitational waves. For a background
$a\propto|\eta|^\alpha$ in conformal time the spectral energy density scales as
\source{GV \S4.3, eq.~(4.105) and Table~2, ll.~4466--4485}
\begin{equation}\label{eq:verdicts:slope}
 \Omega(\omega)\sim\omega^{3-2\nu},\qquad \nu=\left|\alpha-\tfrac12\right| .
\end{equation}
De Sitter inflation ($\alpha=-1$) gives $\nu=3/2$ and a flat spectrum; the dilaton-driven phase has
$\alpha=1/2$, $\nu=0$ and a cubic spectrum $\Omega\sim\omega^3$ \source{GV \S4.3, ll.~4503--4506},
which persists in the low-frequency band of the full model \source{GV \S5.1, eq.~(5.15),
ll.~5040--5053}. The end point is fixed by the curvature at the transition,
$g_1=H_1/M_P\simeq M_s/M_P$ \source{GV \S5.1, eq.~(5.16), ll.~5059--5061}:
\begin{equation}\label{eq:verdicts:end}
 \omega_1\simeq g_1^{1/2}\,10^{11}\,\text{Hz},\qquad \Omega(\omega_1)\simeq10^{-4}g_1^2 .
\end{equation}

\subsection{The programme fixes the string scale}

Candidate C-B assumes (\TC) that scale-factor duality is realized as a pre-big-bang phase and
that the string scale is \eqref{eq:verdicts:ident}, so $M_s\sim1/s$ and
$g_1=\ell_P/s$, i.e.\ $g_1^2=\ell_P^2/s^2=\ell_P H_0/c$. No parameter is left except the position
of a break in the spectrum, which can only lower it below $\omega_1$.

\begin{example}[The numbers of C-B]
$g_1^2=1.616\times10^{-35}/1.3672\times10^{26}=1.18\times10^{-61}$, so $g_1=3.4\times10^{-31}$,
$g_1^{1/2}=5.9\times10^{-16}$, and
\[
 \omega_1\approx5.9\times10^{-5}\,\text{Hz},\qquad \Omega_{\rm GW}(\omega_1)\approx1.2\times10^{-65}.
\]
There are no relic gravitons above $\omega_1$, and below it the spectrum falls like $\omega^3$.
\end{example}

\subsection{The verdict}

C-B was frozen (tag \texttt{stream7-cb-frozen}) before any detector document was pinned, with
the rule: testable iff the prediction reaches the design sensitivity of the most sensitive
detector pinned afterwards. The LISA mission proposal requires the ability to measure
$\Omega=1.3\times10^{-11}(f/10^{-4}\,\text{Hz})^{-1}$ for $0.1\,\text{mHz}<f<2\,\text{mHz}$
\source{LISA, OR7.2, ll.~692--699}, at best $\approx6.5\times10^{-13}$ at $2$\,mHz. The predicted
spectrum ends below LISA's band ($\omega_1<10^{-4}$\,Hz), and its peak lies more than fifty orders
of magnitude below LISA's sensitivity. By the rule, \textbf{C-B is not falsifiable in practice}:
neither confirmed nor excluded.

There is a stronger statement. Read as the Regge slope of the fundamental string,
\eqref{eq:verdicts:ident} puts string resonances at $\hbar c/s\approx4$\,meV, while the CMS dijet
search excludes string resonances below $7.9$\,TeV in the models it tests
\source{\texttt{1911\_03947.txt}, ll.~1044--1048}, a gap of more than $10^{30}$ in $\ap$
(\cref{ch:open}). The premise of C-B is therefore already excluded.

\section{C-C: cosmic F-strings}\label{sec:verdicts:cc}

A fundamental string has tension $\bar\mu=1/(2\pi\ap)$ \source{CMP, eq.~(3.3), l.~505 (with
$p=1$, $q=0$)}, so in units where $G=\ell_P^2$,
\[
 G\mu=\frac{\ell_P^2}{2\pi\ap}=\frac{\ell_P}{2\pi\,c/H_0}
 =\frac{1.616\times10^{-35}}{2\pi\times1.3672\times10^{26}}\approx1.9\times10^{-62}.
\]
The brane-inflation window quoted by Copeland--Myers--Polchinski starts at $G\mu\sim10^{-11}$
\source{CMP \S5.1, eq.~(5.1), ll.~783--784}, and the gravitational-wave reach they quote goes down
to $G\mu\sim10^{-13}$ for LISA \source{CMP, ll.~972--975}. C-C is some fifty orders of magnitude
out of reach, and it shares the excluded premise of C-B.

\section{The common cause, and the geometry itself}\label{sec:verdicts:cause}

All four observables were derived from \eqref{eq:verdicts:ident}. The pattern of the verdicts is
explained by it: taken as a length scale in our universe ($s\approx47\,\mu$m), it conflicts with
torsion-balance and astrophysical bounds (P1); taken with a Hubble-scale infrared cutoff, it
cannot produce acceleration (C-A); taken as the string scale, it puts every stringy signal
below any conceivable sensitivity (C-B, C-C) and is itself excluded by collider searches.
\emph{Every observable the programme can derive today is either excluded or out of reach, and
the cause is the identification $\ap=\ell_P\,c/H_0$.}

A second obstacle is independent of the identification. Type II string theory on $K3\times T^2$
gives $\mathcal N=4$ supersymmetry in four dimensions \source{Asp \S5, l.~2743}, and
$\mathcal N=4$ supersymmetry in four dimensions is not chiral \source{Asp \S5, ll.~2713--2716}.
The chiral fermions of the Standard Model therefore cannot come from this compactification as it
stands (standard; the step from ``the supersymmetry is not chiral'' to ``there is no chiral
matter'' is not spelled out in the pinned lines). A realistic model would need a different
geometry, for instance a Calabi--Yau threefold fibred by K3 surfaces; that is new research.

\section{What the machine has checked}\label{sec:verdicts:lean}

Three modules of \lean{DualScaleCosmology} carry the verdicts. All depend only on
\lean{propext}, \lean{Classical.choice}, \lean{Quot.sound}; each is \TA{} for what it literally
states, which is algebra and arithmetic on stated inputs. None of them states that the
identifications of this chapter are correct.

\subsection{P1}

\lean{DualScaleCosmology/Stream6Verdict.lean}:
\begin{leancode}
noncomputable def p1RadiusSq : ℝ := planckLength_m * hubbleRadius_m

theorem p1_verdict_excluded :
    (30e-6 : ℝ) ^ 2 < p1RadiusSq ∧ (38.6e-6 : ℝ) ^ 2 < p1RadiusSq ∧
      (44e-6 : ℝ) ^ 2 < p1RadiusSq := by ...

theorem p1_factor_needed :
    (30e-6 : ℝ) ^ 2 < (0.64 : ℝ) ^ 2 * p1RadiusSq ∧
      (0.63 : ℝ) ^ 2 * p1RadiusSq < (30e-6 : ℝ) ^ 2 := by ...

theorem p62_selfdual_fixed_point (R s : ℝ) (hR : 0 < R) (hs : 0 < s) :
    R = s ^ 2 / R ↔ R = s := by ...

theorem p62_towers_coincide (s : ℝ) (hs : 0 < s) : 1 / s = s / s ^ 2 := by ...

theorem p62_two_dims_excluded : (1e5 * 1.6e-10 : ℝ) ^ 2 < p1RadiusSq := by ...
\end{leancode}
\begin{leanbox}[Plain-language reading]
\lean{p1\_verdict\_excluded}: the product $\ell_P\cdot c/H_0$ of the two typed-in constants, in
m$^2$, exceeds $(30\,\mu\text{m})^2$, $(38.6\,\mu\text{m})^2$ and $(44\,\mu\text{m})^2$. Proof:
unfold the constants and let \lean{norm\_num} compare decimals.
\lean{p1\_factor\_needed}: a factor $0.64$ still fails the first bound, $0.63$ passes; this
records what a rescue would require, it does not choose it. \lean{p62\_selfdual\_fixed\_point}:
for positive reals, $R=s^2/R$ if and only if $R=s$ (both directions; the proof squares and uses
$(R-s)(R+s)=0$). \lean{p62\_towers\_coincide}: $1/s=s/s^2$. \lean{p62\_two\_dims\_excluded}: the
product exceeds $(10^5\times1.6\times10^{-10}\,\text{m})^2$.
\end{leanbox}

\subsection{C-A}

\lean{DualScaleCosmology/Stream7CA.lean}:
\begin{leancode}
theorem ca_matter_fraction (k H2 rm rL eps : ℝ) (hk : 0 < k)
    (hF : H2 = k * (rm + rL)) (hL : rL = eps * H2 / k) :
    rm = (1 - eps) * H2 / k := by ...

theorem ca_hubble_scaling (k a H2 rm0 eps : ℝ) (hk : 0 < k) (ha : 0 < a)
    (he : eps < 1) (h : rm0 / a ^ 3 = (1 - eps) * H2 / k) :
    H2 = (k * rm0 / (1 - eps)) / a ^ 3 := by ...

theorem ca_deceleration (C a : ℝ) (hC : 0 < C) (ha : 0 < a) :
    -1 - a * deriv (fun x : ℝ => C / x ^ 3) a / (2 * (C / a ^ 3)) = 1 / 2 := by ...

theorem ca_verdict_excluded :
    (1 - DualScaleCosmology.DarkEnergyScale.omegaLambda) / 2
        - DualScaleCosmology.DarkEnergyScale.omegaLambda < 0 ∧
      (1 - DualScaleCosmology.DarkEnergyScale.omegaLambda) / 2
        - DualScaleCosmology.DarkEnergyScale.omegaLambda ≠ 1 / 2 := by ...
\end{leancode}
\begin{leanbox}[Plain-language reading]
\lean{ca\_matter\_fraction} and \lean{ca\_hubble\_scaling} are the two algebraic steps of the
derivation, with $k=8\pi G/3$ left as an arbitrary positive real. \lean{ca\_deceleration} is the
calculus step: with Mathlib's \lean{deriv}, the expression \eqref{eq:verdicts:q} for $E=C/a^3$
equals $1/2$ at every $a>0$; the derivative is computed by the lemma \lean{hasDerivAt\_E}
($\frac{\dd}{\dd a}C/a^3=-3C/a^4$). \lean{ca\_verdict\_excluded}: with $\Omega_\Lambda=0.68885$,
$(1-\Omega_\Lambda)/2-\Omega_\Lambda$ is negative and differs from $1/2$. The Friedmann equation and
the identification \eqref{eq:verdicts:rhoL} are inputs written as hypotheses, not derived.
\end{leanbox}

\subsection{C-B}

\lean{DualScaleCosmology/Stream7CB.lean}:
\begin{leancode}
noncomputable def g1Sq : ℝ := planckLength_m / hubbleRadius_m

noncomputable def slope (α : ℝ) : ℝ := 3 - 2 * |α - 1 / 2|

theorem cb_slope : slope (1 / 2) = 3 ∧ slope (-1) = 0 := by ...

theorem cb_g1_sq_bracket : (1.1e-61 : ℝ) ≤ g1Sq ∧ g1Sq ≤ 1.3e-61 := by ...

theorem cb_peak_below_lisa_band :
    g1Sq * (1e11 : ℝ) ^ 4 < (1e-4 : ℝ) ^ 4 := by ...

noncomputable def omegaPeak : ℝ := 1e-4 * g1Sq

noncomputable def omegaLisaBest : ℝ := 1.3e-11 / 20

theorem cb_not_testable :
    g1Sq * (1e11 : ℝ) ^ 4 < (1e-4 : ℝ) ^ 4 ∧ omegaPeak < omegaLisaBest := by ...
\end{leancode}
\begin{leanbox}[Plain-language reading]
\lean{cb\_slope}: the formula $3-2|\alpha-\tfrac12|$ gives $3$ at $\alpha=1/2$ and $0$ at
$\alpha=-1$ (the formula itself is \TL{} from GV). \lean{cb\_g1\_sq\_bracket}:
$\ell_P/(c/H_0)\in[1.1,1.3]\times10^{-61}$. \lean{cb\_peak\_below\_lisa\_band}: with
$\omega_1=g_1^{1/2}10^{11}$\,Hz, $\omega_1^4<(10^{-4}\,\text{Hz})^4$, i.e.\ $\omega_1<10^{-4}$\,Hz.
\lean{cb\_not\_testable}: in addition, the predicted peak $10^{-4}g_1^2$ is below LISA's best
$1.3\times10^{-11}/20$. The companion theorem \lean{cb\_peak\_amplitude\_unreachable} states the
fifty-orders margin.
\end{leanbox}

\begin{sloppypar}
\noindent C-C uses \lean{CosmicString.}\allowbreak\lean{fStringGmu\_selfDual}
($G\mu=\ell_P/(2\pi L)$ for $\ap=\ell_P L$) and
\lean{CosmicString.}\allowbreak\lean{selfDual\_Gmu\_}\allowbreak\lean{below\_cmp\_window};
the collider exclusion of the premise is
\lean{CosmicString.}\allowbreak\lean{selfDual\_alphaPrime\_}\allowbreak\lean{exceeds\_cms\_ceiling}
(\cref{ch:open}).
\end{sloppypar}

\begin{center}\small
\begin{tabular}{@{}p{0.44\textwidth}c>{\raggedright\arraybackslash}p{0.42\textwidth}@{}}
\toprule
Claim & Tier & Where \\
\midrule
$s^2$ exceeds the three squared bounds; $\kappa<0.64$ needed & \TA & \lean{p1\_verdict\_excluded}, \lean{p1\_factor\_needed} \\
Unique positive fixed point of $R\mapsto s^2/R$ & \TA & \lean{p62\_selfdual\_fixed\_point} \\
$q=1/2$ for $H^2=C/a^3$; $q_0<0$ with $\Omega_\Lambda=0.68885$ & \TA & \lean{ca\_deceleration}, \lean{ca\_verdict\_excluded} \\
$\omega_1<10^{-4}$\,Hz, peak below LISA & \TA & \lean{cb\_not\_testable} \\
Eöt-Wash, neutron-star, LISA, CMS numbers; GV spectrum; CKN bound & \TL & EW, MVV, LISA, CMS, GV, CKN \\
$\ap=\ell_P c/H_0$; $R=s$ as an extra dimension; CKN saturation with $L=c/H$ & \TC & this programme \\
\bottomrule
\end{tabular}
\end{center}

\section{Lessons and open directions}\label{sec:verdicts:lessons}

Four lessons carry over to any theory test. \emph{Freeze before you compute}: the git tag is cheap
and removes every later temptation to adjust. \emph{Declare what you knew}: a retrodiction is
honest if it is called one. \emph{Derive the order-one factors}: here T-duality fixed $\kappa=1$,
and the verdict followed. \emph{Look for the common cause}: four negative verdicts with one shared
input say more about that input than about the four observables.

What would change the situation is new research, not a new convention: dropping or replacing the
identification \eqref{eq:verdicts:ident}, and, for a model of our universe, a compactification
with chiral matter. The mathematics of the earlier chapters (lattices, T-duality, double field
theory, moonshine, dyons) is untouched by these verdicts: it is Tier A where it says so,
whatever happens to the physical hypothesis.

\begin{summarybox}
\begin{itemize}[leftmargin=*,itemsep=1pt]
\item A pre-registered test fixes statement, tests and decision rule by a git tag before the
 verdict; prior knowledge is declared; a changed prediction is a new prediction.
\item Order-one factors must be derived; choosing them after reading a bound is a fit.
\item Without an explicit compactification, the programme yields four observables, all resting on
 $\ap=\ell_P c/H_0$, i.e.\ $s\approx47\,\mu$m.
\item P1 ($s$ as an extra-dimension radius) fails Eöt-Wash and neutron-star bounds and is
 excluded; T-duality fixes $\kappa=1$.
\item C-A (CKN saturation with $L=c/H$) gives $q=1/2$ for every $\varepsilon$; observed $q_0<0$
 excludes it.
\item C-B (pre-big-bang gravitons) ends at $6\times10^{-5}$\,Hz with $\Omega\approx10^{-65}$: not
 falsifiable in practice; C-C gives $G\mu\approx2\times10^{-62}$; both rest on a meV string scale,
 excluded by collider searches.
\item $K3\times T^2$ gives non-chiral $\mathcal N=4$ physics; a realistic model needs another
 geometry.
\end{itemize}
\end{summarybox}

\section*{Exercises}
\addcontentsline{toc}{section}{Exercises}

\begin{exercise}[$\star$]
Recompute $s$ from $\ell_P$ and $c/H_0$ and check the three failures of P1 by comparing squares.
By what factor would $H_0$ have to change for $s$ to fall below $30\,\mu$m? Is such a change
compatible with any measured value of $H_0$ you know? (Answer the second question only after the
first, and say which data you used.)
\end{exercise}

\begin{exercise}[$\star$]
Show that $q=-1-\dot H/H^2$ and that $q=\tfrac12(1+3w)$ for a flat universe dominated by a single
fluid with $p=w\rho$. Use this to confirm that C-A has $w_{\rm eff}=0$.
\end{exercise}

\begin{exercise}[$\star\star$]
Generalize C-A: suppose $\rho_\Lambda=\varepsilon\,3/(8\pi G L^2)$ with $L=c/H$ but the matter
density scales as $a^{-3(1+w_m)}$. Show that $q=\tfrac12(1+3w_m)$ for every $\varepsilon$. What does
this say about any model in which the dark energy is a fixed fraction of the critical density?
\end{exercise}

\begin{exercise}[$\star\star$]
Using \eqref{eq:verdicts:slope}, find the spectral slope for $\alpha=2$ and for $\alpha=-2$, and
classify each background with the table in \source{GV Table~2, ll.~4471--4485}. For $\alpha=1/2$,
how many orders of magnitude does $\Omega$ fall between $\omega_1$ and $10^{-9}$\,Hz (the
pulsar-timing band)?
\end{exercise}

\begin{exercise}[$\star\star$]
Write the frozen document for a \emph{new} prediction P1$'$ in which the extra dimension has
radius $R=\kappa s$ with $\kappa$ derived from some argument of your choice. List the inputs as
measured, derived and assumed, state three tests and a decision rule, and explain which
parts of your derivation would have to be done before the tag.
\end{exercise}

\begin{exercise}[$\star\star$, Lean]
State and prove in Lean 4 with Mathlib: for real $H_2$, $k>0$ and $\varepsilon\ne1$, if
$H_2=k\,\rho_m+\varepsilon H_2$ then $\rho_m=(1-\varepsilon)H_2/k$. Compare your statement with
\lean{ca\_matter\_fraction}: which hypothesis did the module encode differently, and why is the
module's version slightly more general?
\end{exercise}

\begin{exercise}[$\star\star\star$, Lean]
Prove the calculus step of C-A for a general power: for $C>0$, $a>0$ and a natural number $n$,
\begin{leancode}
theorem decel_general (C a : ℝ) (n : ℕ) (hC : 0 < C) (ha : 0 < a) :
    -1 - a * deriv (fun x : ℝ => C / x ^ n) a / (2 * (C / a ^ n))
      = n / 2 - 1 := by sorry
\end{leancode}
and deduce that a universe with $H^2\propto a^{-n}$ accelerates only for $n<2$. Model the proof on
\lean{hasDerivAt\_E}.
\end{exercise}

\begin{exercise}[$\star\star\star$]
C-C and C-B both became untestable because $M_s\sim\hbar c/s$ is tiny. Suppose instead that
$\ap$ were a free parameter. Using the CMP tension formula and \eqref{eq:verdicts:end}, find
the range of $\ap$ for which $G\mu\ge10^{-13}$ and the range for which the pre-big-bang peak
$\Omega(\omega_1)$ exceeds $10^{-12}$. Why would a prediction built this way no longer be a
prediction of the dual-scale programme?
\end{exercise}

\section*{Further reading}
\addcontentsline{toc}{section}{Further reading}
\begin{sloppypar}
\begin{itemize}[leftmargin=*,itemsep=1pt]
\item Torsion-balance bounds on extra dimensions and Yukawa forces: \source{EW, ll.~283--289}.
 Neutron-star bounds and the ``dark dimension'' window: \source{MVV \S3, ll.~318--329, 424--427}.
\item The CKN bound: \source{CKN, ll.~75--80}. Pre-big-bang relic gravitons:
 \source{GV \S4.3, ll.~4466--4506} and \source{GV \S5.1, ll.~5040--5061}.
\item LISA's stochastic-background requirements: \source{LISA, ll.~692--699}. Cosmic F- and
 D-strings: \source{CMP, ll.~497--510, 778--786, 970--976}. String resonances at the LHC:
 \source{\texttt{1911\_03947.txt}, ll.~1044--1048}.
\item $\mathcal N=4$ from $K3\times T^2$ and non-chirality: \source{Asp \S5, ll.~2713--2716,
 2741--2744}.
\item Project documents (in \texttt{docs/}):
 \texttt{STREAM6\_}\allowbreak\texttt{EXPERIMENT\_PLAN.md},
 \texttt{STREAM6\_}\allowbreak\texttt{PREDICTION\_P1.md},
 \texttt{STREAM7\_}\allowbreak\texttt{HYPOTHESIS\_INVENTORY.md},
 \texttt{STREAM7\_}\allowbreak\texttt{PREDICTION\_CA.md},
 \texttt{STREAM7\_}\allowbreak\texttt{PREDICTION\_CB.md},
 \texttt{STREAM67\_ERRATA.md}.
\item Lean sources (in \lean{DualScaleCosmology/}):
 \lean{Stream6Verdict.lean}, \lean{Stream7CA.lean}, \lean{Stream7CB.lean},
 \lean{SelfDualCutoff.lean}, \lean{CosmicString.lean}.
\end{itemize}
\end{sloppypar}
